\documentclass[12pt]{article}

\usepackage[margin=1.25in]{geometry}
\usepackage{setspace}
\usepackage{hyperref}
\usepackage{booktabs}
\usepackage{footnote}
\usepackage[T1]{fontenc}

\hypersetup{
    colorlinks=true,
    linkcolor=blue,
    citecolor=blue,
    urlcolor=blue
}

\doublespacing

\title{George Bancroft's \textit{A Plea for the Constitution of the
United States Wounded in the House of Its Guardians} (1886):\\
The Legal Tender Cases and Their Connection to the\\
Ambiguity of the 5-20 Bonds}
\author{Notes prepared from Bancroft (1886) and related sources}
\date{\today}

\begin{document}

\maketitle
\thispagestyle{empty}

\newpage
\tableofcontents
\newpage

%-----------------------------------------------------------------------
\section{Bancroft: Historian, Statesman, and Monetary Conservative}
%-----------------------------------------------------------------------

George Bancroft (1800--1891) was uniquely positioned to write a
constitutional brief against paper money. As the foremost narrative
historian of the United States---his ten-volume \textit{History of the
United States, from the Discovery of the American Continent} (1834--1874)
was a monument of nineteenth-century scholarship---he brought to monetary
questions not merely the moralism of a hard-money partisan but the
documentary authority of a man who had spent fifty years in the primary
sources of the founding era. He had personally known James Madison, John
Adams, and the Marquis de Lafayette; had served as Secretary of the Navy
under President Polk and as Minister to Prussia and Germany under
Presidents Johnson and Grant; and had been accorded the rare distinction
of floor privileges in the United States Senate by virtue of his name
alone. He was 84 years old when he sat down to write the \textit{Plea},
and he died at 91 having never relented.

The \textit{Plea for the Constitution} (Harper \& Brothers, 1886) is
nominally a response to a single Supreme Court decision---\textit{Juilliard
v.\ Greenman} (1884)---but its ambition is far larger: it is a
comprehensive historical argument that the framers of the Constitution
deliberately, consciously, and by large majorities withheld from the
federal government the power to make paper money legal tender for private
debts, and that a succession of Supreme Court decisions had betrayed this
original design. It is organized in five parts, moving from colonial
monetary history through the Constitutional Convention to the post-Civil
War cases.

%-----------------------------------------------------------------------
\section{The Structure of the Argument}
%-----------------------------------------------------------------------

\subsection{Parts I and II: The Colonial and Revolutionary Experience}

Bancroft opens by tracing the history of American bills of credit from
their first colonial emission in Massachusetts in December 1690---seven
thousand pounds issued to cover the costs of a failed expedition against
Quebec---through the Revolutionary era. The story he tells is one of
inexorable depreciation, repudiation, and moral corruption. Massachusetts
issued its first bills as a war emergency; within nineteen months they
had driven coin from circulation entirely. The pattern repeated in every
colony that followed: bills loaned out on easy terms quickly depreciated,
driving up prices, benefiting debtors at the expense of creditors, and
corrupting commercial morals.

He marshals an impressive series of contemporary witnesses against paper
money. Roger Sherman of Connecticut---one of the framers Bancroft most
admired---wrote in 1752 that money must have certain value, that a
debtor ought not to pay debts with less value than contracted for without
the creditor's consent, and that a fluctuating medium was no better than
false weights and measures. Bancroft invokes Adam Smith's famous judgment
on Pennsylvania's 1723 legal tender law: it ``bears the evident mark of
a scheme of fraudulent debtors to cheat their creditors.'' George
Washington repeatedly condemned paper money in private correspondence,
calling it in 1787 ``foolish and wicked'' and insisting that it is ``by
the substance, not with the shadow of a thing, we are to be benefited.''
Madison characterized it as simultaneously ``unjust,'' ``unconstitutional,''
and ``pernicious.'' The house of delegates of Virginia voted 85--17 in
November 1786 that paper emissions would be ``unjust, impolitic, and
destructive of public and private confidence.''

The revolutionary experience provided the final proof. During the War
of Independence, bills ``became as plenty as oak leaves,'' in John
Adams's phrase; continental dollars were ultimately called in at one
silver dollar for a hundred paper ones---a repudiation so complete that
it served as the decisive empirical exhibit in every subsequent debate.

\subsection{Part III: The Constitutional Convention Closes the Door}

The core of Bancroft's historical argument concerns what happened at
Philadelphia in the summer of 1787. His claim is precise: the framers
did not merely decline to grant the federal government a power to emit
bills of credit as legal tender---they actively and by large majorities
stripped out language that would have permitted it, doing so with full
awareness of what they were doing.

The original draft of the Constitution gave Congress the power ``to
borrow money, and emit bills on the credit of the United States.'' Gouverneur
Morris moved to strike the words ``and emit bills.'' The motion passed
9--2. James Madison noted in his convention journal that the majority
``thought it advisable not to insert this power; thinking it probable
that the omission would not disable the use of public notes so far as
they would be safe and proper, and would only cut off the pretext for a
paper currency, and particularly for making the bills a tender either
for public or private debts.'' Nathaniel Gorham, presiding, asked whether
the deletion would not stop the government from borrowing money on the
credit of the United States; the majority answered that it would not, that
bonds and certificates were plainly permitted---but paper money declared
legal tender for private debts was not.

The states' constitutions told the same story. Article I, Section~10
explicitly forbade states from emitting bills of credit or making anything
but gold and silver coin a tender in payment of debts. This express
prohibition on the states, combined with the congressional power struck
from the draft, was for Bancroft the decisive textual signal: the
Constitution closed the door to paper legal tender, at both levels of
government.

Hamilton, whose monetary nationalism Bancroft might have seemed unlikely
to invoke, had written in 1783---four years before the Convention---that
``to emit an unfunded paper as the sign of value ought not to continue a
formal part of the constitution, nor ever hereafter to be employed; being,
in its nature, pregnant with abuses, and liable to be made the engine of
imposition and fraud.'' That this was the same Hamilton who later built
the First Bank and funded the national debt only strengthened Bancroft's
point: even the most ambitious architects of federal fiscal power drew a
sharp line at paper legal tender.

\subsection{Part IV: The Constitution in the House of Its Guardians}

Part IV is the polemical heart of the work. It traces the Legal Tender
Cases and pronounces them a judicial catastrophe. Bancroft's central
structural device is devastating: he invokes the framers, the earliest
Chief Justices (who were also framers), and the entirety of constitutional
tradition on behalf of Justice Stephen Field's lonely dissent in
\textit{Juilliard}, against the eight-Justice majority.

%-----------------------------------------------------------------------
\section{The Legal Tender Cases: A Trilogy}
%-----------------------------------------------------------------------

The Legal Tender Cases comprise three decisions spanning 1870 to 1884,
each one ratcheting the constitutional permissions of paper money further
than the last.

\subsection{\textit{Hepburn v.\ Griswold} (1870): Chase's Reversal}

The first case in the trilogy is also the most dramatically ironic.
The Legal Tender Act of February~25, 1862---whose drafting Chase had
overseen as Secretary of the Treasury, and which he had initially opposed
before reversing himself under the pressure of war finance---came before
the Supreme Court in \textit{Hepburn v.\ Griswold} (75 U.S.\ 603),
decided February~7, 1870. And the opinion holding the Act unconstitutional
was written by none other than Salmon P.\ Chase, now Chief Justice.

Chase's constitutional argument rested on two pillars. First, the Act
as applied to contracts made \textit{before} its passage impaired the
obligation of those contracts in a manner contrary to the spirit, if
not the letter, of the Constitution. Second, it deprived creditors of
property without due process of law in violation of the Fifth Amendment,
since a creditor who had contracted for payment in gold received
greenbacks worth substantially less. Chase appealed not to any express
constitutional provision but to the ``spirit'' of the instrument---a
methodological move that would later be turned against his own opinion.

The vote was 4--3, and the court that decided it was shorthanded: only
seven justices participated, as vacancies had not been filled. Justice
Miller's dissent pleaded for judicial restraint against Chase's
``abstract and intangible'' spirit-of-the-constitution reasoning. Chase
did not expressly recant his actions as Treasury Secretary; he argued
instead that necessity had compelled what principle now condemned.

Bancroft was entirely sympathetic to Hepburn's result, but the
circumstances of its reversal---fifteen months later, by a reconstituted
court, on a 5--4 vote---gave him as much material for his indictment of
the Supreme Court as the pro-legal-tender decisions themselves.

\subsection{\textit{Knox v.\ Lee} and \textit{Parker v.\ Davis} (1871):
The Court-Packing Reversal}

The reversal came with stunning speed. On the same day that
\textit{Hepburn} was announced (February~7, 1870), President Grant
nominated two new justices to fill existing vacancies: William Strong
of Pennsylvania and Joseph Bradley of New Jersey. When the court reached
its new complement of nine, the Attorney General moved for reargument.
The new court---with Strong and Bradley both participating---reversed
\textit{Hepburn} in \textit{Knox v.\ Lee} and \textit{Parker v.\ Davis}
(79 U.S.\ 457), decided May~1, 1871, by a vote of 5--4.

Justice Strong's majority opinion relied on the implied powers doctrine
of \textit{McCulloch v.\ Maryland}: if Congress has the power to conduct
war and to borrow money on the credit of the United States, and if a
paper currency declared legal tender is an appropriate means of executing
those powers, then the Legal Tender Acts are constitutional under the
Necessary and Proper Clause. Strong addressed the timing coincidence
head-on, chiding Chase for having ``forced'' an earlier decision when the
court was divided and on the verge of receiving new appointees---a
remarkable instance of one majority lecturing a former majority about the
indecorum of the very process by which the current majority came to exist.

Chase dissented, joined by Clifford and Field, in opinions that largely
reiterated the Hepburn majority. Field's dissent---which Bancroft would
later invoke as vindicated by history---was the most constitutionally
principled: the doctrine that whatever is not expressly forbidden may
be exercised would ``change the whole character of our government'' from
one of limited, enumerated powers to one without defined limits in its
relations to the people.

The court-composition question has never been fully resolved. Whether
Grant deliberately packed the court to obtain a reversal of Hepburn, or
whether the nominations were coincidental to the need to fill vacancies,
was debated at the time and remains a matter of controversy. The
structural problem is transparent: on a constitutional question of the
first magnitude, the outcome turned on the personnel of the court rather
than on any change in the law or the facts.

\subsection{\textit{Juilliard v.\ Greenman} (1884): Legal Tender in
Peacetime}

The third and final case in Bancroft's indictment was \textit{Juilliard
v.\ Greenman} (110 U.S.\ 421), decided March~3, 1884. The facts were
prosaic: Juilliard sold one hundred bales of cotton to Greenman for
\$5,100; Greenman attempted to pay in legal tender notes that had
originally been issued during the war, then redeemed and reissued under
the Resumption and Refunding Act of 1878---in peacetime, with no war
emergency to justify them.

Justice Gray's majority opinion, for eight of nine justices, held that
the power of Congress to declare paper money legal tender was not limited
to wartime emergencies but was a permanent attribute of sovereignty. The
key passage, which Bancroft reproduces in his Introduction and treats as
the proximate occasion for his book, reads: ``The power to make the
notes of the government a legal tender in payment of private debts being
one of the powers belonging to sovereignty in other civilized nations,
and not expressly withheld from congress by the constitution; we are
irresistibly impelled to the conclusion that the impressing upon the
treasury notes of the United States the quality of being a legal tender
in payment of private debts is an appropriate means, conducive and plainly
adapted to the execution of the undoubted powers of congress.''

Bancroft's response to this language was contemptuous. The argument from
what ``other civilized nations'' do had no place in the interpretation of
a constitution deliberately designed to withhold powers that every
European monarchy had exercised; the argument from ``not expressly
withheld'' converted a government of enumerated powers into one of
unlimited powers. It was, in his phrase, ``the language of revolution.''

Moreover---and this is the structural irony that runs through Part IV of
the \textit{Plea}---just two years earlier, in October 1882, the very
same nine justices who decided \textit{Juilliard} had unanimously
adopted the Marshallian principle that ``the government of the United
States is one of delegated, limited, and enumerated powers. Therefore
every valid act of congress must find in the Constitution some warrant
for its passage.'' Bancroft remarks acidly that the Juilliard opinion
``is therefore in flagrant antagonism to the constitution of the United
States and to the repeated, continuous, and unanimously given, and
justly respected interpretations of the constitution by the court itself,
under a succession of chiefs without a break.''

Only Justice Field dissented. But Bancroft insists, in a characteristic
rhetorical move, that standing beside Field are ``invincible vouchers for
the rightness of his dissent'': James Wilson, Oliver Ellsworth, and
William Paterson, all three of whom Washington appointed from among the
Convention's framers to serve as the court's earliest interpreters.
Field's dissent, in this framing, is not a minority position but the
transmitted judgment of the founding generation.

\subsection{Part V: What Is To Be Done?}

Bancroft closes by asking whether remedy is possible. His answer is
principally historical and exhortatory: the facts he has assembled are
``beyond the reach of change,'' and the court's only honorable course is
to reconsider. He was under no illusion about the political obstacles.
The \textit{Plea} attracted, in his biographer's words, ``little serious
attention'' at the time of publication. The country had committed itself
to a paper currency; Resumption had been achieved in 1879 and greenbacks
now traded at par with gold; the practical urgency of the constitutional
question seemed to have receded. Bancroft died in 1891 having seen no
reversal.

%-----------------------------------------------------------------------
\subsection{The Pacific Coast Exception: California, Oregon, and the Gold
Contract Doctrine}
%-----------------------------------------------------------------------

While the constitutional drama of \textit{Hepburn}, \textit{Knox}, and
\textit{Juilliard} occupied center stage, a distinct and arguably more
immediately consequential set of developments was unfolding on the Pacific
Coast. California and Oregon resisted the greenback regime with a tenacity
that had no parallel elsewhere in the Union---not primarily through
litigation challenging the constitutionality of the Legal Tender Acts, but
through a combination of state legislation, commercial custom, and judicial
interpretation that rendered those acts largely inoperative for ordinary
contracts within their borders.

\subsubsection*{California's Gold Economy and the Specific Contract Act}

California entered the Union in 1850 as a gold-rush state, and its
commercial life was formed accordingly. No California bank ever suspended
specie payments during the Civil War; gold and silver coin circulated freely
throughout the conflict; and greenbacks, when they crossed the Sierra Nevada
and appeared on Pacific Coast markets, traded at a substantial discount
rather than at par with coin---at the height of wartime inflation, a coined
dollar exchanged for more than two dollars in United States notes.  The
physical proximity of California's economy to the gold and silver mines of
the Sierra Nevada and the Comstock Lode gave it a supply of specie that the
paper-money economy of the East entirely lacked.

The California legislature responded to the Legal Tender Acts by passing,
in 1863, a \textit{Specific Contract Act} that gave legal recognition to the
commercial reality.\footnote{California Statutes 1863, ch.~331.  The act
provided that any contract specifically calling for payment in gold or silver
coin could be ``specifically enforced by execution'' at the creditor's
demand.  Because virtually every commercial instrument in California named
coin as the medium of payment, the act in practice exempted the great
majority of the state's private contracts from the greenback legal tender
provision.}  Under the Act, any contract that specifically called for payment
in gold or silver coin could be specifically enforced in specie: the creditor
could demand coin, not paper, in satisfaction of the debt, and California
courts would compel payment accordingly.  Virtually every commercial
contract executed in California specified gold coin as the medium of
payment---this was the immemorial custom of a gold-producing state---so the
practical effect of the Act was to exempt nearly the whole of ordinary
California commercial practice from the paper-money regime that Congress was
constructing for the rest of the Union.

Murray Rothbard, in his \textit{History of Money and Banking in the United
States} (2002), emphasizes the exceptional character of California's wartime
monetary experience.  California maintained an effective gold standard
throughout the Civil War years, and greenbacks functioned there as a foreign
commodity traded at a discount rather than as lawful circulating money.
State courts consistently refused to apply the legal tender provision to
contracts made and to be performed within California where specie payment had
been expressly stipulated.  Oregon followed a parallel path.  The Oregon
legislature enacted statutes requiring that state taxes be transmitted in
gold and silver coin---``the sheriff shall pay over to the county treasurer
the full amount of the state and school taxes, in gold and silver
coin''---and Oregon courts applied a similar doctrine to coin-specified
commercial contracts.

\subsubsection*{Pre-\textit{Hepburn} Resistance: The California
Supreme Court}

Even before the U.S.\ Supreme Court had addressed the legal tender question,
the California Supreme Court had staked out an important position.  In
\textit{Perry v.\ Washburn} (20 Cal.\ 350, 1862), decided the same year the
Legal Tender Acts were passed, Chief Justice Stephen J.\ Field---who would
later sit on the U.S.\ Supreme Court and write the lone dissent in
\textit{Juilliard}---held that the acts' phrase ``debts, public and private''
referred only to contractual obligations, not to taxes levied by public
authority.  The legal tender provision could not, on Field's construction,
compel a state government or its subdivisions to accept greenbacks in
satisfaction of obligations that were not contractual in origin.  California
had thus ruled the broadest reading of the Legal Tender Acts out of bounds
before the federal courts had said a word on the subject.

\subsubsection*{The Supreme Court Validates the Gold-State Position:
\textit{Lane County}, \textit{Bronson}, and \textit{Butler} (1869)}

The questions raised by California and Oregon state practice reached the
U.S.\ Supreme Court in a cluster of cases decided at the October~1868
term---all of them written by Chief Justice Chase, all of them decided
\textit{before} \textit{Hepburn v.\ Griswold}, and all of them vindicating
the Pacific Coast hard-money position through statutory construction rather
than constitutional adjudication.

In \textit{Lane County v.\ Oregon} (74 U.S.\ 71, 1869), Chase wrote for a
unanimous Court that the federal Legal Tender Acts did not require Oregon to
accept paper money in satisfaction of state tax obligations.  The reasoning
was grounded in federalism: the power of taxation is essential to state
existence; the acts' phrase ``all debts, public and private'' referred to
contractual debts, not to taxes imposed by sovereign authority; and the
express provision making greenbacks receivable for \textit{federal} taxes,
with no corresponding provision for state taxes, implied that Congress had
not intended to override state revenue systems.  Chase cited \textit{Perry
v.\ Washburn} with approval, explicitly endorsing the California Supreme
Court's earlier construction.  States retained full discretion to prescribe
the medium of their own tax collections.

In \textit{Bronson v.\ Rodes} (74 U.S.\ 229, 1869), Chase framed the
gold-contract question with characteristic precision.  A bond made in 1851
explicitly promised payment ``in gold and silver coin, lawful money of the
United States.''  In January 1865, the debtor tendered the nominal amount
due in United States notes; at that moment, one coined dollar was worth two
dollars and a quarter in notes---so the tender was, by Chase's accounting,
worth only six hundred and seventy coined dollars against a debt of fifteen
hundred and seven.  The question was whether the Legal Tender Acts compelled
the creditor to accept such a tender.  Chase held that they did not.  A
contract denominated in coin was a different kind of obligation from a
contract denominated in dollars generally: it was, in legal import, a
contract to deliver a specified weight of standard gold or silver,
``distinguishable in principle'' from a contract to deliver bullion of equal
fineness only by the fact that coins could be counted rather than assayed.
The legal tender laws applied to debts payable in money generally; they did
not apply to obligations that were themselves contracts for specific
commodities.  ``It is not pretended,'' Chase wrote, ``that any real payment
and satisfaction of an obligation to pay fifteen hundred and seven coined
dollars can be made by the tender of paper money worth in the market only
six hundred and seventy coined dollars.''  Critically, Chase refused in
\textit{Bronson} to rule on the constitutionality of the Legal Tender Acts
at all---he reached the pro-gold result through statutory construction,
leaving the constitutional question for another day.

A companion case, \textit{Butler v.\ Horwitz} (74 U.S.\ 258, 1869),
applied the same principle to a pre-Revolutionary ground rent payable in
English golden guineas.  Chase confirmed that the \textit{Bronson} rule
applied to coin-specified contracts made at any date, before or after the
Legal Tender Acts, and that damages on the breach of such contracts should
be assessed in coin rather than in currency.

\subsubsection*{\textit{Trebilcock v.\ Wilson} (1871): The Exception
Survives \textit{Knox}}

The gold-contract doctrine faced a critical test after \textit{Knox v.\ Lee}
had reversed \textit{Hepburn} and given legal tender full constitutional
sanction.  In \textit{Trebilcock v.\ Wilson} (79 U.S.\ 687, 1871)---decided
in the same term as \textit{Knox}, with Justice Field writing for a
majority---the Court held unanimously that a promissory note payable ``in
specie'' was a coin contract fully within the rule of \textit{Bronson}.  The
term ``in specie'' was merely descriptive of the kind of dollars
promised---coined rather than paper---and a contract for coined dollars
could not be discharged by tender of notes, whatever the constitutional
status of those notes as legal tender for debts payable in money generally.
\textit{Knox} had validated the legal tender power in the abstract; it had
not disturbed \textit{Bronson}'s distinction between commodity contracts and
money contracts.  The Court in \textit{Trebilcock} ruled expressly that the
Legal Tender Acts applied only ``to debts which are payable in money
generally, and not to obligations payable in commodities or obligations of
any other kind,'' and that coin contracts fell decisively in the latter
category.  Justice Bradley dissented, arguing that all contracts payable in
``money of the United States'' in whatever description were subject to the
legal tender provision; Justices Miller and Bradley together were the
minority.

\subsubsection*{Chase's Strategy and the Field Connection}

The Pacific Coast episode illuminates Chase's jurisprudential strategy in
the legal tender controversy as a whole.  The three Chase opinions of
1869---\textit{Lane County}, \textit{Bronson}, and \textit{Butler}---were
handed down before \textit{Hepburn} was decided.  In each, Chase reached
a hard-money outcome through statutory construction, deliberately declining
to address the constitutional question.  He created a doctrinal
alternative---the gold-contract exception---that insulated Pacific Coast
creditors from greenback inflation without requiring the Court to strike
down an act of Congress during wartime.  Only in February~1870, after the
war was over and Resumption was under discussion, did he bring the
constitutional question forward directly in \textit{Hepburn}.  One reading
of the sequence is that Chase was constructing, brick by brick, a
jurisprudential structure that would minimize the scope of the legal tender
power while avoiding open confrontation with Congress on the constitutional
question until the political moment was more propitious.

Equally notable is the career arc of Stephen J.\ Field through these
cases.  As Chief Justice of the California Supreme Court, Field wrote the
\textit{Perry v.\ Washburn} opinion in 1862, holding that ``all debts,
public and private'' did not include state tax obligations---the precise
construction later adopted by Chase in \textit{Lane County}.  Appointed to
the U.S.\ Supreme Court in 1863, Field became the Court's most consistent
defender of hard-money rights: the dissenter in \textit{Knox v.\ Lee}, the
lone dissenter in \textit{Juilliard}, and---in the interval---the author of
the \textit{Trebilcock} majority that preserved the \textit{Bronson}
exception even after \textit{Knox} had settled the constitutional question
adversely.  The gold-contract doctrine was in this sense Field's most
durable achievement in a jurisprudential career devoted to limiting the
paper-money power; it outlasted every other element of the hard-money
constitutional program.

\subsubsection*{Connection to Bancroft's Argument}

Bancroft does not discuss the California and Oregon episode at length in the
\textit{Plea}---his argument is focused on the pure constitutional question
of congressional power, not the alternative doctrinal route of coin
contracts---but the Pacific Coast cases are closely related to his central
theme.  The gold-contract doctrine represents a partial victory for the
hard-money position: it preserved, for those creditors who had expressly
bargained for coin, the specific performance that the Constitution (in
Bancroft's view) guaranteed to all creditors as a matter of original design.
That the victory was partial---that it protected the sophisticated drafter
who had written ``gold and silver coin'' into the instrument but not the
ordinary counterparty who had contracted for ``dollars'' without further
specification---was precisely the fragility that Bancroft's broad
constitutional argument was intended to cure.  The doctrine required
creditors to have anticipated, years or decades in advance, the introduction
of a depreciated paper currency and to have protected themselves by express
contract language; it offered no remedy to those who had trusted that the
legal framework would not be changed in a manner prejudicial to their
existing rights.

The Pacific Coast exception thus stands, in the larger history of the Legal
Tender Cases, as a durable but narrow survival of the hard-money principle:
durable because it survived \textit{Knox}, \textit{Juilliard}, and the
entire subsequent development of legal tender jurisprudence (the
\textit{Bronson} doctrine remained operative until Congress abrogated
gold-payment clauses in contracts by joint resolution in June~1933); narrow
because it applied only to explicit coin contracts, leaving the great
majority of dollar-denominated obligations subject to the paper-money regime
against which Bancroft had argued in vain.  From the perspective of Hall and
Sargent (2014), the 5-20 bonds themselves were coin contracts in the
\textit{Bronson} sense for the interest component---the statute expressly
required coin for interest---but not for the principal, which was
denominated in dollars without further specification.  The \textit{Bronson}
exception thus protected 5-20 holders on the interest side automatically,
while leaving the principal ambiguity to be resolved by the political
process that produced the Public Credit Act of 1869.  On the Pacific Coast,
the Specific Contract Act and the gold-contract doctrine together created
the legal framework within which that resolution was, for most ordinary
commercial obligations, entirely unnecessary: coin was specified, coin was
received, and the paper-money question never arose.

%-----------------------------------------------------------------------
\section{The Chase Irony and Its Significance}
%-----------------------------------------------------------------------

The figure of Salmon P.\ Chase runs through the entire episode like a
thread of constitutional tragedy. As Secretary of the Treasury he
overcame his initial legal scruples and signed the Legal Tender Act of
February~25, 1862 into policy, marketing the 5-20 bonds to the public
with the implicit promise---never guaranteed by statute---that principal
as well as interest would be paid in gold. As Chief Justice he voted
in \textit{Hepburn} to strike down the very Act he had championed,
finding that it impaired private contracts against the spirit of the
Constitution. And then he was reversed by a court reconstituted through
Grant's appointments, his \textit{Hepburn} majority dissolving fifteen
months after it was formed.

Bancroft does not dwell on the irony with personal relish---he treats
Chase's conduct as symptomatic of a constitutional culture that had
lost its moorings rather than as individual hypocrisy. But he clearly
regards the chase episode as illustrating his general thesis: that when
great powers are exercised without a clear constitutional warrant, relying
instead on expedience and executive reputation, the legal and moral
foundations crumble at the first serious challenge. Chase's reputational
commitment to gold repayment of the 5-20 principal, and his
constitutional commitment to the invalidity of legal tender paper, were
both extra-statutory---both existed as assurances he gave without
explicit legislative authority---and both were undone by the same
political logic that had originally forced him to accept the Legal Tender
Act as a war measure.

%-----------------------------------------------------------------------
\section{Connection to the 5-20 Bond Ambiguity}
%-----------------------------------------------------------------------

Reading Bancroft's \textit{Plea} alongside the history of the 5-20 bonds
illuminates a structural theme that runs across both: the creation of
deliberate ambiguity as a short-term political solution, and the
downstream costs of resolving it.

\subsection{Two Ambiguities in Parallel}

The 5-20 bonds created one well-documented ambiguity: the Act of
February~25, 1862 specified that interest would be paid in coin, but was
entirely silent on principal. This was almost certainly an intentional
legislative omission---a compromise that bought the votes needed for
passage while preserving fiscal flexibility in a war with uncertain
outcome. Chase then bridged the statutory gap with reputational
commitment, marketing the bonds as coin-redeemable throughout.

The Legal Tender Acts created a parallel constitutional ambiguity: whether
the federal government's power to declare paper legal tender was a
permanent attribute of sovereignty (as \textit{Juilliard} ultimately
held) or a temporary emergency measure only (as Chase's position in
\textit{Hepburn} implied). Congress in 1862 did not address this question
because the urgency of war finance crowded out constitutional deliberation.
The same dynamic that produced the 5-20 principal silence---competing
coalitions, unresolvable disagreement, and the pressure of necessity---
produced the legal tender ambiguity.

\subsection{Resolution by Political Force Rather Than Constitutional Logic}

Both ambiguities were resolved not by the logic of the original texts but
by political outcomes that were not guaranteed in advance. The 5-20
principal ambiguity was resolved by Grant's election in November 1868
and the Public Credit Act of March~18, 1869. The legal tender
constitutional ambiguity was resolved by Grant's court appointments,
producing the \textit{Knox v.\ Lee} reversal in 1871, and definitively
settled in \textit{Juilliard} in 1884.

In both cases, the ``resolution'' favored the party with the stronger
electoral coalition---hard-money Republicans in the former, a broad
congressional majority unwilling to constrain its own fiscal powers
in the latter. Bancroft's constitutional argument is, at bottom, a claim
that this is not how constitutional settlements should work: the framers
had already resolved the paper money question at Philadelphia, and the
Supreme Court's obligation was to honor that resolution rather than to
defer to contemporary political majorities.

\subsection{Commitment, Ambiguity, and the Chase Paradox}

There is a deeper theoretical connection between the two ambiguities that
Bancroft does not articulate explicitly but that his narrative makes
visible. In both cases, reputational commitment was used as a substitute
for statutory or constitutional commitment---and in both cases the
substitute proved fragile when tested.

In the 5-20 episode, Chase's reputation as a hard-money Democrat served
as the bonding mechanism for the implicit gold-repayment promise. The
commitment was credible so long as Chase remained Secretary and the
political context held---but it was not self-enforcing, as Pendleton's
``Ohio Idea'' immediately demonstrated. In the legal tender episode,
the constitutional ``commitment'' to limited federal power over money
was credible so long as the Court composition remained stable---but it
dissolved within fifteen months when two vacancies were filled.

Bancroft's point can be restated in the terms of modern commitment theory:
constitutional provisions that are clear and explicit are hard constraints;
provisions that rest on implication or the ``spirit'' of the document are
soft constraints that yield to political pressure. The 5-20 silence was
a soft constraint on the principal repayment question; the omission of
``emit bills'' from Article I, Section~8 was a soft constraint on the
federal paper money question. Both were resolved by political outcomes.
What Bancroft hoped for---and what he believed the Convention had
actually achieved---was a hard constraint: a textual prohibition clear
enough to resist exactly the kind of ``irresistible'' reasoning by
which Justice Gray found paper money to be a legitimate exercise of
sovereignty.

\subsection{The Widened Scope of the Problem}

From the perspective of Hall and Sargent (2014), the 5-20 bond ambiguity
was a relatively contained fiscal discrimination problem: it concerned
who would receive gold versus greenbacks in the settlement of a specific
class of government obligations, and it was resolved within seven years
of the relevant bond issuance. The legal tender constitutional ambiguity
was a far more general problem: it concerned whether the government had
the power to expand the class of obligations it could discharge in
paper, indefinitely and in perpetuity.

\textit{Juilliard}'s holding---that the legal tender power is an
attribute of national sovereignty not restricted to war---meant that the
fiscal discriminations documented by Hall and Sargent became potentially
permanent features of the constitutional order rather than emergency
expedients. Bancroft saw this clearly. His fear was not that the
particular bondholders of 1862 would be cheated (they had been made
whole, with interest in coin and principal effectively in gold by the
Public Credit Act); his fear was that the Court had legitimized an
instrument of fiscal discretion that could be deployed against any class
of future creditors without limit.

The resolution of the 5-20 ambiguity in 1869 thus sat uneasily alongside
the \textit{Juilliard} ruling of 1884. Grant's Public Credit Act had
pledged gold repayment for the 5-20s specifically; but \textit{Juilliard}
said that Congress had the power to do precisely what Pendleton had
proposed---pay public debts in legal tender notes---whenever a future
political majority chose to do so. The statutory commitment of 1869 and
the constitutional permission of 1884 pointed in opposite directions.
Bancroft's \textit{Plea} is, among other things, an argument that the
1884 decision retroactively undermined the constitutional basis for the
honor that Grant and the Republicans had claimed in 1869.

%-----------------------------------------------------------------------
\section{Bancroft's Intellectual Legacy on These Questions}
%-----------------------------------------------------------------------

Bancroft's \textit{Plea} was not well received in its day, and the
constitutional position it advocated has never been revived by the Supreme
Court. \textit{Juilliard} has been repeatedly reaffirmed; the legal tender
power of Congress is not seriously contested in contemporary constitutional
law. The Federal Reserve System, created in 1913, operates entirely on
the legal foundation that Bancroft believed the framers had forbidden.

Yet the \textit{Plea} retains its interest for several reasons. As
historical scholarship, it is a remarkable synthesis of colonial monetary
history that anticipated the professional historiography by decades; the
basic narrative of colonial paper money depreciation that Bancroft
assembled in 1886 remains largely accurate. As constitutional argument,
it is the fullest statement of the ``original design'' case against paper
legal tender, and its documentation of the Philadelphia debates over
the ``emit bills'' clause provides the primary evidence still cited by
those who argue (as a minority of constitutional scholars continues to do)
that the framers' intent was stricter than subsequent jurisprudence has
acknowledged.

But perhaps its deepest interest is as a case study in the long-run
consequences of short-term ambiguity. The Civil War Congress of 1862
resolved its immediate fiscal emergency by passing the Legal Tender Act
without confronting the constitutional question clearly; the Supreme Court
of 1870--1884 resolved the resulting constitutional ambiguity in a manner
that Bancroft believed---with considerable historical evidence on his
side---was inconsistent with the founders' design. And the original
motivation for the 1862 ambiguity---war finance under conditions of
political uncertainty---was precisely what Bancroft's historical narrative
had catalogued as the universal context in which paper money had always
been adopted and always been abused, from Massachusetts Bay in 1690
to continental dollars in 1780.

The 5-20 bonds and the Legal Tender Cases are thus two aspects of a
single underlying problem in American fiscal constitutionalism: the
difficulty of making credible commitments about money when the
institutional mechanisms for doing so---clear statutory language, explicit
constitutional prohibition, and a politically insulated judiciary---
are themselves products of the political process whose behavior they
are designed to constrain.

%-----------------------------------------------------------------------
\section*{References}
%-----------------------------------------------------------------------

\begin{description}
\setlength\itemsep{0.5em}

\item Bancroft, George. 1886. \textit{A Plea for the Constitution of the
United States of America Wounded in the House of Its Guardians}. New
York: Harper \& Brothers. Digitized copy at Internet Archive:
\texttt{archive.org/details/apleaforconstit00bancgoog}. Full text also
at \texttt{constitution.org/2-Authors/gb/gb-plea.htm}.

\item Chase, Salmon P. 1870. Majority opinion in \textit{Hepburn
v.\ Griswold}, 75 U.S.\ (8 Wall.) 603.

\item Field, Stephen J. 1871. Dissenting opinion in \textit{Knox v.\
Lee} and \textit{Parker v.\ Davis}, 79 U.S.\ (12 Wall.) 457.

\item Grant, Ulysses S. 1869. ``Public Credit Act,'' March 18, 1869,
41st Cong., 1st Sess., Chapter 1, 16 Stat. 1.

\item California. Statutes of California, 1863, Chapter~331. Specific
Contract Act (``An Act to authorize the specific enforcement of certain
contracts'').

\item Chase, Salmon P. 1869. Majority opinions in \textit{Lane County
v.\ Oregon}, 74 U.S.\ (7 Wall.) 71; \textit{Bronson v.\ Rodes}, 74 U.S.\
(7 Wall.) 229; and \textit{Butler v.\ Horwitz}, 74 U.S.\ (7 Wall.) 258.
All three decided at the October~1868 term of the Supreme Court.

\item Field, Stephen J. 1862. Majority opinion as Chief Justice of
California in \textit{Perry v.\ Washburn}, 20 Cal.\ 350 (California Supreme
Court, 1862).

\item Field, Stephen J. 1871. Majority opinion in \textit{Trebilcock
v.\ Wilson}, 79 U.S.\ (12 Wall.) 687, decided at the October~1870 term of
the Supreme Court.

\item Gray, Horace. 1884. Majority opinion in \textit{Juilliard
v.\ Greenman}, 110 U.S.\ 421 (March~3, 1884).

\item Hall, George J. and Thomas J. Sargent. 2014. ``Fiscal
Discriminations in Three Wars.'' \textit{Journal of Monetary Economics}
61: 148--166.

\item Lowenstein, Roger. 2022. \textit{Ways and Means: Lincoln and His
Cabinet and the Financing of the Civil War}. New York: Penguin Press.

\item McCulloch v. Maryland, 17 U.S.\ (4 Wheat.) 316 (1819). Marshall,
C.J.

\item Rothbard, Murray N. 2002. \textit{A History of Money and Banking in
the United States: The Colonial Era to World War~II}. Auburn, AL: Ludwig
von Mises Institute.

\item Sherman, Roger. 1752. \textit{A Caveat Against Injustice, or an
Inquiry into the Evil Consequences of a Fluctuating Medium of Exchange}.
New Haven. Reprinted in Bancroft (1886).

\item Strong, William. 1871. Majority opinion in \textit{Knox v.\ Lee}
and \textit{Parker v.\ Davis}, 79 U.S.\ (12 Wall.) 457.

\item United States. Congress. \textit{Congressional Globe}, 37th
Congress, 2nd Session (1862). Washington, D.C.: Library of Congress
American Memory.

\item United States. Statutes at Large, Vol.\ XII, Chap.\ XXXIII
(February 25, 1862). Legal Tender Act. St.\ Louis Fed FRASER:
\texttt{fraser.stlouisfed.org/title/legal-tender-act-1107}.

\end{description}

\end{document}
